\documentclass[a4paper,11pt]{article}

\usepackage[margin=2cm]{geometry}
\usepackage[italian]{babel}
\usepackage[utf8]{inputenc}
\usepackage[T1]{fontenc}
\usepackage{lmodern}

\usepackage{enumerate}

\usepackage{amsmath}
\usepackage{amsfonts}
\usepackage{amssymb}
\usepackage{amsthm}

\usepackage[dvipsnames]{xcolor}
\usepackage{listings}
\lstdefinestyle{mystyle}{
    commentstyle=\color{ForestGreen},
    keywordstyle=\color{magenta},
    numberstyle=\tiny\color{gray},
    stringstyle=\color{purple},
    basicstyle=\ttfamily\footnotesize,
    breakatwhitespace=false,
    breaklines=true,
    captionpos=b,
    keepspaces=true,
    numbers=left,
    numbersep=5pt,
    showspaces=false,
    showstringspaces=false,
    showtabs=false,
    tabsize=2
}
\lstset{style=mystyle}
\lstset{framextopmargin=50pt,frame=bottomline}

\begin{document}

\section{Problem}\label{sec:orig}

The health care department has to hire new doctors. In the department there are
m specializations (orthopedics, pediatrics, etc.). For each specialization $j$ a
minimum number of new doctors $b_j$ are required. There are $n$ candidate
doctors. Each doctor $i$ can cover several specializations (at the same time).
$S_i$ denotes the set of specializations that doctor $i$ covers. The salary cost
of each doctor $i$, if hired, is $c_i$.

The problem of the health care department manager is to hire a subset of doctors
that guarantees to cover all requirements at minimum cost.

\subsection{Formulation}\label{sec:orig-formulation}

To represent the problem we use a bipartite graph. One side of the graph will be
the set of doctors $D$ (indexed by $i$), while on the other the set of all
specializations $S$ (indexed by $j$). The set of arcs between the two sets will
be $A$, which will be defined in terms of the parameter $S_i$ as such:

\[
  A = \{ (i,j) | i\in D \land j\in S_i \}
\]

Thus the resulting graph will be $G = (D\times S, A)$.

As parameters we will use the aforementioned $S_i$ representing the set of
specialization of doctor $i$, the salaries of each doctor $c_i$ and the minimum
hires of each specialization $b_j$.

As variables we will use two binary variables:

\begin{enumerate}
  \item $x_{ij}$ indicates whether the edge $(i,j)$ has been chosen, this
    translates to whether the doctor $i$ has been chosen to cover specialty $j$;
  \item $y_i$ indicates whether a doctor has been hired to work for any
    specialty
\end{enumerate}

The constraints and objective are thus as follows:

\begin{align}
  \min \sum_{i\in D} c_i y_i & \\
  \sum_{(i,j) \in \textit{BS}(j)} x_{ij} & \geq b_j & \forall j \in S \label{eqn:doctor-req} \\
  y_i - \sum_{j\in S} x_{ij} & \geq 0 & \forall i \in D \label{eqn:linking} \\
  x_{ij} \in \{0,1\}, &\quad y_i \in \{0,1\}
\end{align}

Where $\textit{BS}(j)$ is the backwards star of node $j$.
Equation~\eqref{eqn:doctor-req} represents the doctor requirement constraint while
equation~\eqref{eqn:linking} the linking constraint for $y_i$.

\subsection{Some python code}

Some python code has been write and can be run in a python notebook using
\texttt{mip}.

\begin{lstlisting}[language=Python]
import networkx as nx
import matplotlib.pyplot as plt
import numpy as np
import mip

%matplotlib inline

np.random.seed(4321)
n = 5
m = 3

### Sets and parameters
D = [i for i in range(n)]
S = [j for j in range(m)]

Si = {i: set(np.random.randint(0, m, size=3)) for i in D}
C = [np.random.randint(1, 10) for i in D]
A = [(i, j) for i in D for j in Si[i]]

B = np.random.randint(1, max([len(Si[i]) for i in Si]), size=m)

m = mip.Model()
### Variables
vars = {(i, j): m.add_var(var_type=mip.BINARY) for (i, j) in A} # x_ij
selected = [m.add_var(var_type=mip.BINARY) for i in D]          # y_i

### Constraints
for j in S:
    m.add_constr(mip.xsum(vars[i, j] for (i, j1) in A if j == j1) >= B[j])
for i in D:
    m.add_constr(selected[i] >= mip.xsum(vars[i,j] for (i1, j) in A if i1 == i))

### Objective
m.objective = mip.minimize(mip.xsum(C[i] * selected[i] for i in D))
m.optimize()
\end{lstlisting}

\section{Problem variant}

Different from what has been outlined in~\ref{sec:orig}, we can now leave some
specialization uncovered and activate interregional contracts. For each
specialization the cost of the contract if $t_j$. In case the contract is
activated, there is no need to hire any doctor to cover $b_j$ positions. Our
objective now is to minimize the total cost of both new hires and contracts.

\subsection{Formulation}

We can keep most of what we built in~\ref{sec:orig-formulation}. However, we
need to add some ingredients in light of the new requirements:

\begin{description}
  \item[Parameters] We add a new parameter $t_j$ indicating the cost of the
    contract for specialization $j$.
  \item[Variables] We add a new binary variable $z_j$ indicating whether a
    contract for specialization $j$ has been chosen.
\end{description}

We can the rewrite the objective and constraints as follows:

\begin{align}
  \min \sum_{i\in D} c_i y_i & + \sum_{j\in S} t_j z_j  \\
  \sum_{(i,j) \in \textit{BS}(j)} x_{ij} & \geq b_j (1 - z_j) & \forall j \in S \label{eqn:doctor-req-var} \\
  y_i - \sum_{j\in S} x_{ij} & \geq 0 & \forall i \in D \label{eqn:linking-var} \\
  x_{ij} \in \{0,1\}, &\quad y_i \in \{0,1\}, \quad z_j \in \{0,1\}
\end{align}

Where~\eqref{eqn:doctor-req-var} is the new doctor requirement constraint
while~\eqref{eqn:linking-var} is the same linking constraint as before.

\subsection{Some other python code}

Like before, some more python code is provided so the model can be tested inside
a python notebook. It is simply a modified version of the previous listing.

\begin{lstlisting}[language=Python]
import networkx as nx
import matplotlib.pyplot as plt
import numpy as np
import mip

%matplotlib inline

np.random.seed(4321)
n = 5
m = 3

### Sets and parameters
D = [i for i in range(n)]
S = [j for j in range(m)]

Si = {i: set(np.random.randint(0, m, size=3)) for i in D}
C = [np.random.randint(1, 10) for i in D]
T = [np.random.randint(5, 15) for j in S]

A = [(i,j) for i in D for j in Si[i]]

B = np.random.randint(1, max([len(Si[i]) for i in Si]), size=m)

m = mip.Model()
### Variables
vars = {(i,j): m.add_var(var_type=mip.BINARY) for (i,j) in A} # x_ij
selected = [m.add_var(var_type=mip.BINARY) for i in D]        # y_i
contract = [m.add_var(var_type=mip.BINARY) for j in S]        # z_j

### Constraints
for j in S:
    m.add_constr(
        mip.xsum(vars[i, j] for (i, j1) in A if j == j1) >= B[j]*(1 - contract[j]))
for i in D:
    m.add_constr(selected[i] >= mip.xsum(vars[i,j] for (i1,j) in A if i1 == i))

### Objective
m.objective = mip.minimize(
    mip.xsum(C[i]*selected[i] for i in D) + mip.xsum(T[j]*contract[j] for j in S))
m.optimize()
\end{lstlisting}

\end{document}
